\scanpage{291}
\[
b)\quad\varphi_N:\mathcal M_N\longrightarrow\mathcal M
\]
\[
(10.2)\quad\left\{\begin{array}{l}
(\varphi_N(Q))_i=\operatorname{Hom}^i_{\mathcal M}(\tau_{N-i}(S),Q\inserted{i_NQ})
\bigl(\longrightarrow\operatorname{Hom}^i_{\mathcal M}(S,Q)=Q_i\bigr)\\
\hspace{4em}=\operatorname{Hom}^{\unsure{b|0}}_{\mathcal M}(\tau_{N-i}(S)[-i],Q)
\end{array}\right.
\]
Avec \deleted{\unsure{leur}} structure de module provenant (pour $k\in\mathbb Z$, $u\in S_k^{\unsure{b|h}}$) des opérateurs ``multiplication par $u$'' :
\diagram{\begin{tikzcd}\tau_{N-i}(S)&\tau_{N-i-k}(S)\arrow[l,"{\text{deg }k}"]\end{tikzcd}}
en transposant; on a un hom (\unsure{dans} \deleted{\unsure{mod}} $\mathcal M$ décrit dans (10.2)
\[
(10.2.1)\qquad\varphi_N(Q)\longrightarrow i_N(Q)
\]
d'où en appliquant $\tau_N$
\[
(10.2.2)\qquad\tau_N\varphi_NQ\underset{\text{adj.}}{\overset{\sim}{\longrightarrow}}Q
\]
qui est un isom (c'est un hom d'\unsure{adjoint})

L'autre hom d'adjonction
\diagram{\begin{tikzcd}
\deleted{\varphi_N\tau_NP\longleftarrow}&P\arrow[r]&\varphi_N\tau_NP
\end{tikzcd}}
est \unsure{déduit|défini} des \deleted{\unsure{l'inclusion}} \inserted{hom}
\[
\deleted{\tau_NP\longrightarrow P}\qquad\operatorname{Hom}^i_{\mathcal M}(\tau_{N-\deleted{\unsure{i}}}S,P)
\]
(10.2.3)
\diagram{\begin{tikzcd}[column sep=small]
P_i\arrow[r]\arrow[dr,"\unsure{\sim|2}"']&(\varphi_N\tau_NP)_i\simeq\operatorname{Hom}^i(\tau_{N-i}S,\tau_NP)\\
&\operatorname{Hom}^i(S,P)\arrow[u]
\end{tikzcd}}
\margin{tout hom de degré $i$ \deleted{de $S$ dans $P$} $\tau_{N-i}S$ dans $P$ est à valeurs dans $\tau_{N\deleted{\unsure{-i}}}P\subset P$}

c) (NB On n'utilisera pas $q_N$, on le met par acquit de conscience) :
\[
q_N(P)=P/\text{sous-module engendré par les }P_i\text{ pour }i<N
\]
\margin{292}
